\begin{definition} \label{an1:def:open-cover}
  A family $(G_i)_{i\in I}$ is an \textbf{open cover} of $E$ if each $G_i$ is open and $E\subseteq\bigcup_{i\in I}G_i$.
  For $S\subseteq I$, the subfamily indexed by $S$ is a \textbf{subcover} if $E\subseteq\bigcup_{i\in S}G_i$.
  The family retains its original indices; finite index sets are represented by \texttt{Finset}.
\end{definition}

\begin{definition} \label{an1:def:compact}
  A set $E$ is \textbf{compact} if every open cover admits a finite subcover.
  We use \texttt{IsCompact E}. Its finite-open-subcover characterization is recorded below, in the Lean context \texttt{universe u} and \texttt{variable \{X : Type u\} [MetricSpace X]}.
  Index types in the same universe suffice for this characterization; the standard predicate also handles covers indexed in other universes.
\end{definition}

% Lean source: Textbooks.MathematicalAnalysis1.Chapter01.CompactSets.isCompact_iff_finite_subcover
\begin{lemma} \label{an1:lem:isCompact-iff-finite-subcover}
  $E$ is compact if and only if every family of open sets covering it has a finite subfamily still covering it.

  \begin{lean}
    theorem isCompact_iff_finite_subcover (E : Set X) :
        IsCompact E ↔ ∀ {ι : Type u} (U : ι → Set X),
          (∀ i, IsOpen (U i)) → E ⊆ ⋃ i, U i → ∃ s : Finset ι, E ⊆ ⋃ i ∈ s, U i
  \end{lean}
\end{lemma}

\begin{proof}
  The finite-subcover interface expresses compactness with the same open sets, coverage inclusion, and finite set of indices in both directions.

  \begin{lean}
    constructor

    · intro h ι U hU hc
      exact h.elim_finite_subcover U hU hc

    · intro h
      apply isCompact_of_finite_subcover
      exact h
  \end{lean}
\end{proof}

% Lean source: Textbooks.MathematicalAnalysis1.Chapter01.CompactSets.compact_subspace
\begin{proposition} \label{an1:thm:compactness-preserved-under-subspace}
  For $E\subseteq Y\subseteq X$, $E$ is compact in the subspace $Y$ if and only if it is compact in $X$.

  \begin{lean}
    theorem compact_subspace (E Y : Set X) (hEY : E ⊆ Y) :
        IsCompact ((Subtype.val : Y → X) ⁻¹' E) ↔ IsCompact E
  \end{lean}
\end{proposition}

\begin{proof}
  First suppose $E$ is compact in $Y$. Restrict an ambient open cover to $Y$. An ambient ball contained in a cover member restricts to a subspace ball, so each restricted member is open.

  \begin{lean}
    classical

    constructor

    · intro h
      apply isCompact_of_finite_subcover
      intro ι U hU hc

      have ho : ∀ i, IsOpen ((Subtype.val : Y → X) ⁻¹' U i) := by
        intro i
        apply Metric.isOpen_iff.mpr
        intro p hp

        obtain ⟨r, hr, hs⟩ := Metric.isOpen_iff.mp (hU i) p.val hp

        exact ⟨r, hr, fun q hq => hs hq⟩
  \end{lean}

  The restricted family covers $E$ in $Y$. A finite subcover there gives the same finite collection of ambient members covering $E$ in $X$.

  \begin{lean}
    have hcov : (Subtype.val : Y → X) ⁻¹' E ⊆ ⋃ i, Subtype.val ⁻¹' U i := by
      intro p hp

      obtain ⟨i, hi⟩ := mem_iUnion.mp (hc hp)

      exact mem_iUnion.mpr ⟨i, hi⟩

    obtain ⟨s, hs⟩ := h.elim_finite_subcover _ ho hcov

    refine ⟨s, ?_⟩
    intro p hp

    obtain ⟨i, hi⟩ := mem_iUnion.mp (hs (show (⟨p, hEY hp⟩ : Y) ∈ Subtype.val ⁻¹' E from hp))
    obtain ⟨his, hpi⟩ := mem_iUnion.mp hi

    exact mem_iUnion.mpr ⟨i, mem_iUnion.mpr ⟨his, hpi⟩⟩
  \end{lean}

  Conversely, let a family of subspace-open sets cover $E$. Each member is the intersection with $Y$ of an ambient open set, by the relative-open-set proposition.

  \begin{lean}
    · intro h
      apply isCompact_of_finite_subcover
      intro ι U hU hc

      have hex : ∀ i, ∃ V : Set X, IsOpen V ∧ Subtype.val '' U i = V ∩ Y := by
        intro i
        apply (relative_open _ Y (by
          rintro x ⟨p, _, rfl⟩
          exact p.property)).mp

        have heq : (Subtype.val : Y → X) ⁻¹' (Subtype.val '' U i) = U i :=
          preimage_image_eq _ Subtype.val_injective

        rw [heq]
        exact hU i
  \end{lean}

  Choose those ambient open sets. They cover $E$ in $X$, so compactness gives finitely many of them.

  \begin{lean}
    choose V hV heq using hex

    have hcov : E ⊆ ⋃ i, V i := by
      intro p hp

      obtain ⟨i, hi⟩ := mem_iUnion.mp (hc (show (⟨p, hEY hp⟩ : Y) ∈ Subtype.val ⁻¹' E from hp))

      have hm : p ∈ Subtype.val '' U i := ⟨⟨p, hEY hp⟩, hi, rfl⟩

      rw [heq i] at hm
      exact mem_iUnion.mpr ⟨i, hm.1⟩
  \end{lean}

  For a point of $E\subseteq Y$, membership in a chosen ambient set is membership in its intersection with $Y$. Thus the corresponding finitely many original members cover $E$ in the subspace.

  \begin{lean}
    obtain ⟨s, hs⟩ := h.elim_finite_subcover V hV hcov

    refine ⟨s, ?_⟩
    intro p hp

    obtain ⟨i, hi⟩ := mem_iUnion.mp (hs hp)
    obtain ⟨his, hpi⟩ := mem_iUnion.mp hi

    have hm : p.val ∈ Subtype.val '' U i := (heq i).symm ▸ ⟨hpi, p.property⟩

    obtain ⟨q, hq, he⟩ := hm

    have hqp : q = p := Subtype.ext he

    exact mem_iUnion.mpr ⟨i, mem_iUnion.mpr ⟨his, hqp ▸ hq⟩⟩
  \end{lean}

\end{proof}

% Lean source: Textbooks.MathematicalAnalysis1.Chapter01.CompactSets.compact_closed
\begin{proposition} \label{an1:thm:compact-set-is-closed}
  Every compact subset $E$ of a metric space is closed.

  \begin{lean}
    theorem compact_closed (E : Set X) (hE : IsCompact E) : IsClosed E
  \end{lean}
\end{proposition}

\begin{proof}
  To prove that $E^{c}$ is open, fix $p\notin E$. For each $q\in E$, let $r_q=d(q,p)/2>0$. The balls $B(q,r_q)$ cover $E$.

  \begin{lean}
    classical

    apply isOpen_compl_iff.mp
    apply Metric.isOpen_iff.mpr
    intro p hp

    let r : E → ℝ := fun q => dist q.val p / 2

    have hr : ∀ q : E, 0 < r q := by
      intro q
      exact half_pos (dist_pos.mpr (fun heq => hp (heq ▸ q.property)))

    have hcover : E ⊆ ⋃ q : E, ball q.val (r q) := by
      intro q hq
      exact mem_iUnion.mpr ⟨⟨q, hq⟩, by simpa only [mem_ball, dist_self] using hr ⟨q, hq⟩⟩
  \end{lean}

  Take a finite subcover and a common positive lower bound $\delta$ for its radii. This is also possible when the finite subcover is empty.

  \begin{lean}
    obtain ⟨s, hs⟩ := hE.elim_finite_subcover _ (fun q => ball_open _ _) hcover
    obtain ⟨δ, hδ, hle⟩ := finite_positive_lower_bound s r (fun q _ => hr q)

    refine ⟨δ, hδ, ?_⟩
  \end{lean}

  If $x\in B(p,\delta)\cap E$, one of the selected balls contains $x$. Then $d(q,p)\leq d(q,x)+d(x,p)<r_q+\delta\leq 2r_q=d(q,p)$, a contradiction. Hence $B(p,\delta)\subseteq E^{c}$.

  \begin{lean}
    intro x hx hxe

    obtain ⟨q, hq⟩ := mem_iUnion.mp (hs hxe)
    obtain ⟨hqs, hxq⟩ := mem_iUnion.mp hq

    have hd := dist_triangle q.val x p

    rw [dist_comm q.val x] at hd

    have hδq := hle q hqs

    change dist x p < δ at hx
    change dist x q.val < dist q.val p / 2 at hxq
    change δ ≤ dist q.val p / 2 at hδq
    linarith
  \end{lean}

\end{proof}

% Lean source: Textbooks.MathematicalAnalysis1.Chapter01.CompactSets.compact_bounded
\begin{proposition} \label{an1:thm:compact-set-is-bounded}
  Every compact subset of a metric space is bounded, with the stated empty-space convention.

  \begin{lean}
    theorem compact_bounded (E : Set X) (hE : IsCompact E) : Bornology.IsBounded E
  \end{lean}
\end{proposition}

\begin{proof}
  If $X$ is nonempty, choose $p\in X$. The balls $B(p,n+1)$, for $n\in\N$, cover $E$: every real distance is below some natural number by the Archimedean property.

  \begin{lean}
    classical

    apply (isBounded_iff E).mpr
    intro hX

    obtain ⟨p⟩ := hX

    have hc : E ⊆ ⋃ n : ℕ, ball p (n + 1 : ℝ) := by
      intro q _

      obtain ⟨n, hn⟩ := exists_nat_gt (dist q p)

      exact mem_iUnion.mpr ⟨n, by
        change dist q p < (n : ℝ) + 1
        linarith⟩
  \end{lean}

  Choose a finite subcover and set $R=1+\sum_{n\in S}(n+1)$. This is positive, even when $S$ is empty.

  \begin{lean}
    obtain ⟨s, hs⟩ := hE.elim_finite_subcover _ (fun n => ball_open _ _) hc

    let R : ℝ := 1 + ∑ n ∈ s, (n + 1 : ℝ)

    have hnonneg : 0 ≤ ∑ n ∈ s, (n + 1 : ℝ) :=
      Finset.sum_nonneg (fun n _ => by positivity)

    refine ⟨p, R, by
      dsimp [R]
      linarith, ?_⟩
  \end{lean}

  Each $q\in E$ belongs to a selected ball. Its radius is at most the sum of all selected radii, so $d(q,p)<R$. Thus $E\subseteq B(p,R)$.

  \begin{lean}
    intro q hq

    obtain ⟨n, hn⟩ := mem_iUnion.mp (hs hq)
    obtain ⟨hns, hdist⟩ := mem_iUnion.mp hn

    have hle : (n + 1 : ℝ) ≤ ∑ k ∈ s, (k + 1 : ℝ) :=
      Finset.single_le_sum (f := fun k : ℕ => (k + 1 : ℝ)) (fun k _ => by positivity) hns

    change dist q p < R
    change dist q p < (n + 1 : ℝ) at hdist
    dsimp [R]
    linarith
  \end{lean}

\end{proof}

% Lean source: Textbooks.MathematicalAnalysis1.Chapter01.CompactSets.closed_subset_compact
\begin{proposition} \label{an1:thm:closed-subset-of-compact-set-is-compact}
  If $E$ is compact, $F$ is closed, and $F\subseteq E$, then $F$ is compact.

  \begin{lean}
    theorem closed_subset_compact (E F : Set X) (hE : IsCompact E)
        (hF : IsClosed F) (hFE : F ⊆ E) : IsCompact F
  \end{lean}
\end{proposition}

\begin{proof}
  Take an open cover of $F$. Add the open set $F^{c}$, giving the added member its own index distinct from the original indices.

  \begin{lean}
    classical

    apply isCompact_of_finite_subcover
    intro ι U hU hcover

    let V : Option ι → Set X := fun i => match i with
      | none => Fᶜ
      | some j => U j

    have hV : ∀ i, IsOpen (V i) := by
      intro i
      cases i with
      | none => exact hF.isOpen_compl
      | some j => exact hU j
  \end{lean}

  The enlarged family covers $E$: points in $F$ are covered by the original family, and the others lie in $F^{c}$. Select a finite subcover of $E$.

  \begin{lean}
    have hc : E ⊆ ⋃ i, V i := by
      intro x _
      by_cases hx : x ∈ F

      · obtain ⟨i, hi⟩ := mem_iUnion.mp (hcover hx)

        exact mem_iUnion.mpr ⟨some i, hi⟩

      · exact mem_iUnion.mpr ⟨none, hx⟩

    obtain ⟨s, hs⟩ := hE.elim_finite_subcover V hV hc
  \end{lean}

  Discard the added member and keep the selected original indices. A point of $F$ cannot lie in $F^{c}$, so it remains covered. This gives a finite subcover of $F$.

  \begin{lean}
    refine ⟨s.biUnion Option.toFinset, ?_⟩
    intro x hx

    obtain ⟨i, hi⟩ := mem_iUnion.mp (hs (hFE hx))
    obtain ⟨his, hxi⟩ := mem_iUnion.mp hi
    cases i with
    | none => exact False.elim (hxi hx)
    | some j =>
      exact mem_iUnion.mpr ⟨j, mem_iUnion.mpr
        ⟨Finset.mem_biUnion.mpr
          ⟨some j, his, by simp only [Option.toFinset_some, Finset.mem_singleton]⟩, hxi⟩⟩
  \end{lean}

\end{proof}

% Lean source: Textbooks.MathematicalAnalysis1.Chapter01.CompactSets.infinite_subset_limit_point
\begin{theorem} \label{an1:thm:infinite-subset-of-compact-set-has-limit-point}
  If $E$ is compact and $F\subseteq E$ is infinite, then $F^{\prime}\cap E\ne\varnothing$.

  \begin{lean}
    theorem infinite_subset_limit_point (E F : Set X) (hE : IsCompact E)
        (hFE : F ⊆ E) (hF : F.Infinite) : (derivedSet F ∩ E).Nonempty
  \end{lean}
\end{theorem}

\begin{proof}
  Assume there is no such limit point. Then each $p\in E$ has a radius $r_p>0$ for which the only possible point of $F$ inside $B(p,r_p)$ is $p$ itself. Otherwise every radius would have a distinct witness, making $p$ a limit point.

  \begin{lean}
    classical

    by_contra hn

    have hlocal : ∀ p : E, ∃ r > 0, ∀ q ∈ F, dist q p.val < r → q = p.val := by
      intro p
      by_contra h
      push Not at h
      apply hn
      refine ⟨p.val, ?_, p.property⟩
      rw [mem_derivedSet_iff]
      intro r hr

      obtain ⟨q, hq, hd, hne⟩ := h r hr

      exact ⟨q, hq, hne, hd⟩
  \end{lean}

  These balls cover $E$, so choose finitely many of them.

  \begin{lean}
    choose r hr hsingle using hlocal

    have hc : E ⊆ ⋃ p : E, ball p.val (r p) := by
      intro p hp
      exact mem_iUnion.mpr ⟨⟨p, hp⟩, by simpa only [mem_ball, dist_self] using hr ⟨p, hp⟩⟩

    obtain ⟨s, hs⟩ := hE.elim_finite_subcover _ (fun p => ball_open _ _) hc
  \end{lean}

  Every point of $F$ lies in one of the selected balls and must equal its center. Hence $F$ is contained in the finite set of selected centers, contradicting its infinitude.

  \begin{lean}
    apply hF
    apply (s.finite_toSet.image Subtype.val).subset
    intro q hq

    obtain ⟨p, hp⟩ := mem_iUnion.mp (hs (hFE hq))
    obtain ⟨hps, hd⟩ := mem_iUnion.mp hp

    exact ⟨p, hps, (hsingle p q hq hd).symm⟩
  \end{lean}

\end{proof}

% Lean source: Textbooks.MathematicalAnalysis1.Chapter01.CompactSets.compact_finite_intersection
\begin{theorem} \label{an1:thm:finite-compact-intersection-property}
  Suppose $(K_i)_{i\in I}$ are compact and every finite intersection $\bigcap_{i\in S}K_i$ is nonempty, including the intersection over no indices. Then $\bigcap_{i\in I}K_i$ is nonempty.

  \begin{lean}
    theorem compact_finite_intersection {ι : Type*} (K : ι → Set X)
        (hK : ∀ i, IsCompact (K i))
        (hfinite : ∀ s : Finset ι, (⋂ i ∈ s, K i).Nonempty) :
        (⋂ i, K i).Nonempty
  \end{lean}
\end{theorem}

\begin{proof}
  First suppose $I$ has an index $i_0$. If the full intersection were empty, the complements of the $K_i$ would cover $K_{i_0}$. These complements are open because compact sets are closed.

  \begin{lean}
    classical

    by_cases hi : Nonempty ι

    · obtain ⟨i₀⟩ := hi
      by_contra hn

      have hc : K i₀ ⊆ ⋃ i, (K i)ᶜ := by
        intro p _
        by_contra hp
        apply hn
        refine ⟨p, mem_iInter.mpr ?_⟩
        intro i
        by_contra hpi
        exact hp (mem_iUnion.mpr ⟨i, hpi⟩)
  \end{lean}

  A finite subcover would make the intersection over its indices together with $i_0$ empty. The finite-intersection hypothesis instead supplies a point in all these sets, contradicting that subcover.

  \begin{lean}
    obtain ⟨s, hs⟩ := (hK i₀).elim_finite_subcover _
      (fun i => (compact_closed _ (hK i)).isOpen_compl) hc

    obtain ⟨p, hp⟩ := hfinite (insert i₀ s)

    have hp0 := mem_iInter.mp (mem_iInter.mp hp i₀) (Finset.mem_insert_self _ _)

    obtain ⟨i, hi⟩ := mem_iUnion.mp (hs hp0)
    obtain ⟨his, hpi⟩ := mem_iUnion.mp hi

    exact hpi (mem_iInter.mp (mem_iInter.mp hp i) (Finset.mem_insert_of_mem his))
  \end{lean}

  If $I$ is empty, the hypothesis for the empty finite index set supplies a point of $X$. It automatically belongs to every member of the empty family.

  \begin{lean}
    · obtain ⟨p, _⟩ := hfinite ∅

      refine ⟨p, mem_iInter.mpr ?_⟩
      intro i
      exact False.elim (hi ⟨i⟩)
  \end{lean}

\end{proof}

% Lean source: Textbooks.MathematicalAnalysis1.Chapter01.CompactSets.nested_intervals
\begin{lemma} \label{an1:lem:nested-intervals}
  Let $(a_n)$ and $(b_n)$ be real sequences indexed by $n\in\N$ with $a_n\leq a_{n+1}\leq b_{n+1}\leq b_n$. There exists $x\in\R$ with $a_n\leq x\leq b_n$ for every $n$, so the nested intervals have nonempty intersection.

  \begin{lean}
    theorem nested_intervals (a b : ℕ → ℝ)
        (h : ∀ n, a n ≤ a (n + 1) ∧ a (n + 1) ≤ b (n + 1) ∧ b (n + 1) ≤ b n) :
        ∃ x : ℝ, ∀ n, a n ≤ x ∧ x ≤ b n
  \end{lean}
\end{lemma}

\begin{proof}
  The left endpoints increase and the right endpoints decrease. For any $n,m$, compare the indices: if $n\leq m$, then $a_n\leq a_m\leq b_m$; otherwise $a_n\leq b_n\leq b_m$. Thus each $b_m$ bounds all left endpoints.

  \begin{lean}
    have ha : Monotone a := monotone_nat_of_le_succ (fun n => (h n).1)

    have hb : Antitone b := antitone_nat_of_succ_le (fun n => (h n).2.2)

    have hab : ∀ n, a n ≤ b n := fun n => (h n).1.trans ((h n).2.1.trans (h n).2.2)

    have hcross : ∀ n m, a n ≤ b m := by
      intro n m
      rcases le_total n m with hnm | hmn

      · exact (ha hnm).trans (hab m)
      · exact (hab n).trans (hb hmn)
  \end{lean}

  The set of left endpoints contains $a_0$ and is bounded above by $b_0$. Its supremum exists. Each $a_n$ is below that supremum, and each $b_n$ is an upper bound for it, giving a point in every interval.

  \begin{lean}
    have hne : (range a).Nonempty := ⟨a 0, mem_range_self 0⟩

    have hbd : BddAbove (range a) := ⟨b 0, by
      rintro x ⟨n, rfl⟩
      exact hcross n 0⟩

    refine ⟨sSup (range a), ?_⟩
    intro n
    exact ⟨le_csSup hbd (mem_range_self n), csSup_le hne (by
      rintro x ⟨m, rfl⟩
      exact hcross m n)⟩
  \end{lean}

\end{proof}

% Lean source: Textbooks.MathematicalAnalysis1.Chapter01.CompactSets.nested_cells
\begin{theorem} \label{an1:thm:nested-k-cells}
  Let $k\in\N$ and let real endpoints satisfy $a_{n,j}\leq a_{n+1,j}\leq b_{n+1,j}\leq b_{n,j}$ for every $n\in\N$ and $0\leq j<k$. Then the cells $\prod_{j=0}^{k-1}[a_{n,j},b_{n,j}]$ have nonempty intersection.

  \begin{lean}
    theorem nested_cells (k : ℕ) (a b : ℕ → Fin k → ℝ)
        (h : ∀ n j, a n j ≤ a (n + 1) j ∧
          a (n + 1) j ≤ b (n + 1) j ∧ b (n + 1) j ≤ b n j) :
        ∃ x : EuclideanSpace ℝ (Fin k), ∀ n j, a n j ≤ x j ∧ x j ≤ b n j
  \end{lean}
\end{theorem}

\begin{proof}
  For each coordinate, apply the nested-interval result and choose a point belonging to every interval in that coordinate. The resulting vector belongs to every cell. When $k=0$, the choice has no coordinates and produces the unique empty tuple.

  \begin{lean}
    have hex : ∀ j, ∃ x : ℝ, ∀ n, a n j ≤ x ∧ x ≤ b n j :=
      fun j => nested_intervals (fun n => a n j) (fun n => b n j) (fun n => h n j)

    choose x hx using hex

    exact ⟨WithLp.toLp 2 x, fun n j => hx j n⟩
  \end{lean}

\end{proof}

We now prove compactness of cells from open covers. We first prove the interval case using a supremum, then combine finitely many coordinates. This also accounts for degenerate intervals and the zero-dimensional cell.
The next supporting argument applies more generally to topological spaces: their open sets contain the empty and whole sets and are closed under arbitrary unions and finite intersections. Metric spaces have these properties by the preceding results. Compactness still means the finite-open-subcover property. A map is continuous in the open-set sense: inverse images of open sets are open. Only this defining property is needed.

% Lean source: Textbooks.MathematicalAnalysis1.Chapter01.CompactSets.compact_image
\begin{lemma} \label{an1:lem:compact-image}
  For topological spaces $A$ and $B$, if $K\subseteq A$ is compact and $f:A\to B$ is continuous, then $f(K)$ is compact.

  \begin{lean}
    theorem compact_image {A B : Type*} [TopologicalSpace A] [TopologicalSpace B]
        (K : Set A) (hK : IsCompact K) (f : A → B) (hf : Continuous f) :
        IsCompact (f '' K)
  \end{lean}
\end{lemma}

\begin{proof}
  Pull an open cover of $f(K)$ back along $f$. The inverse images are open and cover $K$, so finitely many cover it.

  \begin{lean}
    apply isCompact_of_finite_subcover
    intro ι U hU hc

    have hcov : K ⊆ ⋃ i, f ⁻¹' U i := by
      intro x hx

      obtain ⟨i, hi⟩ := mem_iUnion.mp (hc ⟨x, hx, rfl⟩)

      exact mem_iUnion.mpr ⟨i, hi⟩

    obtain ⟨s, hs⟩ := hK.elim_finite_subcover _ (fun i => (hU i).preimage hf) hcov
  \end{lean}

  Every point of $f(K)$ has a preimage in $K$. That preimage belongs to a selected inverse image, so the point belongs to the corresponding selected cover member.

  \begin{lean}
    refine ⟨s, ?_⟩
    rintro y ⟨x, hx, rfl⟩

    obtain ⟨i, hi⟩ := mem_iUnion.mp (hs hx)
    obtain ⟨his, hxi⟩ := mem_iUnion.mp hi

    exact mem_iUnion.mpr ⟨i, mem_iUnion.mpr ⟨his, hxi⟩⟩
  \end{lean}

\end{proof}

% Lean source: Textbooks.MathematicalAnalysis1.Chapter01.CompactSets.singleton_compact
\begin{lemma} \label{an1:lem:singleton-compact}
  A singleton in any metric space is compact.

  \begin{lean}
    theorem singleton_compact {A : Type*} [MetricSpace A] (a : A) :
        IsCompact ({a} : Set A)
  \end{lean}
\end{lemma}

\begin{proof}
  One member of any cover contains its sole point. The subfamily containing only that member already covers the singleton.

  \begin{lean}
    apply isCompact_of_finite_subcover
    intro ι U _ hc

    obtain ⟨i, hi⟩ := mem_iUnion.mp (hc (mem_singleton a))

    refine ⟨{i}, ?_⟩
    intro x hx
    rw [mem_singleton_iff] at hx
    subst x
    exact mem_iUnion.mpr ⟨i, mem_iUnion.mpr ⟨Finset.mem_singleton_self _, hi⟩⟩
  \end{lean}

\end{proof}

% Lean source: Textbooks.MathematicalAnalysis1.Chapter01.CompactSets.interval_compact
\begin{lemma} \label{an1:lem:interval-compact}
  If $a,b\in\R$ and $a\leq b$, then $[a,b]$ is compact.

  \begin{lean}
    theorem interval_compact (a b : ℝ) (hab : a ≤ b) : IsCompact (Icc a b)
  \end{lean}
\end{lemma}

\begin{proof}
  Fix an open cover. Let $T$ consist of the endpoints $x\in[a,b]$ for which $[a,x]$ has a finite subcover. A cover member containing $a$ covers $[a,a]$, so $a\in T$.

  \begin{lean}
    classical

    apply isCompact_of_finite_subcover
    intro ι U hU hc

    let T : Set ℝ := {x | x ∈ Icc a b ∧ ∃ s : Finset ι, Icc a x ⊆ ⋃ i ∈ s, U i}

    obtain ⟨i₀, hi₀⟩ := mem_iUnion.mp (hc (show a ∈ Icc a b from ⟨le_rfl, hab⟩))

    have ha : a ∈ T := by
      refine ⟨⟨le_rfl, hab⟩, {i₀}, ?_⟩
      intro x hx

      have heq : x = a := le_antisymm hx.2 hx.1

      subst x
      exact mem_iUnion.mpr ⟨i₀, mem_iUnion.mpr ⟨Finset.mem_singleton_self _, hi₀⟩⟩
  \end{lean}

  The set $T$ is nonempty and bounded above by $b$. Let $c=\sup T$; then $a\leq c\leq b$. Choose a cover member containing $c$ and a radius $r>0$ whose ball about $c$ stays in that member.

  \begin{lean}
    have hne : T.Nonempty := ⟨a, ha⟩

    have hb : BddAbove T := ⟨b, fun x hx => hx.1.2⟩

    let c := sSup T

    have hac : a ≤ c := le_csSup hb ha

    have hcb : c ≤ b := csSup_le hne (fun x hx => hx.1.2)

    obtain ⟨i, hi⟩ := mem_iUnion.mp (hc (show c ∈ Icc a b from ⟨hac, hcb⟩))
    obtain ⟨r, hr, hball⟩ := Metric.isOpen_iff.mp (hU i) c hi
  \end{lean}

  Choose $t\in T$ with $c-r/2<t$. Set $d=\min(b,c+r/2)$, which is at least $c$. We will extend a finite cover of $[a,t]$ to $[a,d]$ by adjoining the chosen member.

  \begin{lean}
    obtain ⟨t, ht, hct⟩ := exists_lt_of_lt_csSup hne (show c - r / 2 < sSup T by
      dsimp [c]
      linarith)

    obtain ⟨s, hs⟩ := ht.2

    let d := min b (c + r / 2)

    have hcd : c ≤ d := le_min hcb (by linarith)

    have hd : d ∈ T := by
      refine ⟨⟨hac.trans hcd, min_le_left _ _⟩, insert i s, ?_⟩
  \end{lean}

  A point $x\in[a,d]$ with $x\leq t$ is covered by the old subcover. Otherwise $c-r/2<t<x\leq c+r/2$, so $x$ is in the chosen ball. Therefore $d\in T$.

  \begin{lean}
    intro x hx
    by_cases hxt : x ≤ t

    · obtain ⟨j, hj⟩ := mem_iUnion.mp (hs ⟨hx.1, hxt⟩)

      obtain ⟨hjs, hxj⟩ := mem_iUnion.mp hj

      exact mem_iUnion.mpr ⟨j, mem_iUnion.mpr ⟨Finset.mem_insert_of_mem hjs, hxj⟩⟩

    · have hxu : x ≤ c + r / 2 := hx.2.trans (min_le_right _ _)

      have hxball : x ∈ ball c r := by
        rw [mem_ball, Real.dist_eq, abs_lt]
        constructor <;> linarith

      exact mem_iUnion.mpr ⟨i, mem_iUnion.mpr ⟨Finset.mem_insert_self _ _, hball hxball⟩⟩
  \end{lean}

  Since $c$ is an upper bound of $T$, we have $d\leq c$. If $c<b$, both $b$ and $c+r/2$ are greater than $c$, which would give $d>c$. Hence $c=b$ and $d=b$. The finite subcover witnessing $d\in T$ covers the whole interval.

  \begin{lean}
    have hdc : d ≤ c := le_csSup hb hd

    have hbc : b ≤ c := by
      by_contra h

      have hlt : c < d := lt_min (lt_of_not_ge h) (by linarith)

      exact (not_lt_of_ge hdc) hlt

    have hdb : d = b := le_antisymm (min_le_left _ _) (hbc.trans hcd)

    rw [hdb] at hd
    exact hd.2
  \end{lean}

\end{proof}

% Lean source: Textbooks.MathematicalAnalysis1.Chapter01.CompactSets.compact_product
\begin{lemma} \label{an1:lem:compact-product}
  For compact sets $K\subseteq A$ and $L\subseteq B$ in metric spaces, $K\times L$ is compact in the product metric $d((x,y),(x^{\prime},y^{\prime}))=\max(d(x,x^{\prime}),d(y,y^{\prime}))$.

  \begin{lean}
    theorem compact_product {A B : Type*} [MetricSpace A] [MetricSpace B]
        (K : Set A) (L : Set B) (hK : IsCompact K) (hL : IsCompact L) :
        IsCompact (K ×ˢ L)
  \end{lean}
\end{lemma}

\begin{proof}
  Fix an open cover and $x\in K$. For each $y\in L$, choose a cover member and a product ball of radius $r_y>0$ about $(x,y)$ contained in it.

  \begin{lean}
    classical

    apply isCompact_of_finite_subcover
    intro ι U hU hc

    have htube : ∀ x : K, ∃ δ > 0, ∃ s : Finset ι,
        ∀ x' ∈ ball x.val δ, ∀ y ∈ L, (x', y) ∈ ⋃ i ∈ s, U i := by
      intro x

      have hlocal : ∀ y : L, ∃ i : ι, ∃ r > 0, ball (x.val, y.val) r ⊆ U i := by
        intro y

        obtain ⟨i, hi⟩ := mem_iUnion.mp
          (hc (show (x.val, y.val) ∈ K ×ˢ L from ⟨x.property, y.property⟩))
        obtain ⟨r, hr, hs⟩ := Metric.isOpen_iff.mp (hU i) (x.val, y.val) hi

        exact ⟨i, r, hr, hs⟩
  \end{lean}

  The balls about the $y$ coordinates cover $L$. Compactness selects finitely many; take a common positive lower bound $\delta$ for their radii.

  \begin{lean}
    choose i r hr hs using hlocal

    have hcov : L ⊆ ⋃ y : L, ball y.val (r y) := by
      intro y hy
      exact mem_iUnion.mpr ⟨⟨y, hy⟩, by simpa only [mem_ball, dist_self] using hr ⟨y, hy⟩⟩

    obtain ⟨t, ht⟩ := hL.elim_finite_subcover _ (fun y => ball_open _ _) hcov
    obtain ⟨δ, hδ, hle⟩ := finite_positive_lower_bound t r (fun y _ => hr y)
  \end{lean}

  If $d(x^{\prime},x)<\delta$ and $y^{\prime}\in L$, choose a selected ball containing $y^{\prime}$. Both coordinate distances are below its radius, so $(x^{\prime},y^{\prime})$ lies in its original cover member. Thus finitely many members cover the whole strip $B(x,\delta)\times L$.

  \begin{lean}
    refine ⟨δ, hδ, t.image i, ?_⟩
    intro x' hx' y hy

    obtain ⟨z, hz⟩ := mem_iUnion.mp (ht hy)
    obtain ⟨hzt, hyz⟩ := mem_iUnion.mp hz

    have hpair : (x', y) ∈ ball (x.val, z.val) (r z) := by
      change max (dist x' x.val) (dist y z.val) < r z
      exact max_lt (hx'.trans_le (hle z hzt)) hyz

    exact mem_iUnion.mpr ⟨i z, mem_iUnion.mpr ⟨Finset.mem_image.mpr ⟨z, hzt, rfl⟩, hs z hpair⟩⟩
  \end{lean}

  Choose such a strip and finite index set for each $x\in K$. Their balls cover $K$, so finitely many strips suffice.

  \begin{lean}
    choose δ hδ s hs using htube

    have hcov : K ⊆ ⋃ x : K, ball x.val (δ x) := by
      intro x hx
      exact mem_iUnion.mpr ⟨⟨x, hx⟩, by simpa only [mem_ball, dist_self] using hδ ⟨x, hx⟩⟩

    obtain ⟨t, ht⟩ := hK.elim_finite_subcover _ (fun x => ball_open _ _) hcov
  \end{lean}

  Take the union of the finite index sets for those strips. For any $(x,y)\in K\times L$, one selected strip contains it, and its associated members cover it. The resulting finite family covers the product.

  \begin{lean}
    refine ⟨t.biUnion s, ?_⟩
    rintro ⟨x, y⟩ ⟨hx, hy⟩

    obtain ⟨z, hz⟩ := mem_iUnion.mp (ht hx)
    obtain ⟨hzt, hxz⟩ := mem_iUnion.mp hz

    obtain ⟨i, hi⟩ := mem_iUnion.mp (hs z x hxz y hy)
    obtain ⟨his, hxyi⟩ := mem_iUnion.mp hi

    exact mem_iUnion.mpr ⟨i, mem_iUnion.mpr ⟨Finset.mem_biUnion.mpr ⟨z, hzt, his⟩, hxyi⟩⟩
  \end{lean}

\end{proof}

% Lean source: Textbooks.MathematicalAnalysis1.Chapter01.CompactSets.pi_cell_compact
\begin{lemma} \label{an1:lem:pi-cell-compact}
  For $n\in\N$ and real endpoints $a_j\leq b_j$ for $0\leq j<n$, the coordinate box $\{x: a_j\leq x_j\leq b_j\text{ for all }j\}$ is compact in the finite-coordinate product metric.

  \begin{lean}
    theorem pi_cell_compact (n : ℕ) (a b : Fin n → ℝ) (hab : ∀ j, a j ≤ b j) :
        IsCompact {x : Fin n → ℝ | ∀ j, a j ≤ x j ∧ x j ≤ b j}
  \end{lean}
\end{lemma}

\begin{proof}
  Induct on $n$. With no coordinates, every vector is the same empty tuple, so the box is a compact singleton.

  \begin{lean}
    induction n with
    | zero =>
      have heq : {x : Fin 0 → ℝ | ∀ j, a j ≤ x j ∧ x j ≤ b j} = {fun j => Fin.elim0 j} := by
        ext x
        simp only [mem_ofPred_eq, mem_singleton_iff]
        constructor

        · intro _
          funext j
          exact Fin.elim0 j

        · intro _ j
          exact Fin.elim0 j

      rw [heq]
      exact singleton_compact _
  \end{lean}

  For $n+1$ coordinates, separate coordinate $0$ from the remaining $n$. The map adjoining that first coordinate is continuous: each output-coordinate distance is bounded by the maximum product distance of the inputs, so $\delta=\varepsilon$ suffices. The interval for coordinate $0$ and the remaining box are compact, so their product is compact.

  \begin{lean}
    | succ n ih =>
      let K := {x : Fin n → ℝ | ∀ j, a j.succ ≤ x j ∧ x j ≤ b j.succ}

      let f : ℝ × (Fin n → ℝ) → (Fin (n + 1) → ℝ) := fun p => Fin.cons p.1 p.2

      have hf : Continuous f := by
        apply Metric.continuous_iff.mpr
        intro p ε hε
        refine ⟨ε, hε, ?_⟩
        intro q hq
        apply (dist_pi_lt_iff hε).mpr
        intro j

        have hpair : max (dist q.1 p.1) (dist q.2 p.2) < ε := hq

        refine Fin.cases ?_ (fun k => ?_) j

        · exact (max_lt_iff.mp hpair).1
        · exact (dist_le_pi_dist q.2 p.2 k).trans_lt (max_lt_iff.mp hpair).2

      have hk := compact_product (Icc (a 0) (b 0)) K (interval_compact _ _ (hab 0))
        (ih (fun j => a j.succ) (fun j => b j.succ) (fun j => hab j.succ))
  \end{lean}

  The image of this product is exactly the desired box: its first coordinate lies in the first interval, and the tail lies in the remaining intervals. Conversely, split any vector in the box into its first coordinate and tail. Compactness of continuous images completes the induction.

  \begin{lean}
    have heq : f '' (Icc (a 0) (b 0) ×ˢ K) =
        {x : Fin (n + 1) → ℝ | ∀ j, a j ≤ x j ∧ x j ≤ b j} := by
      ext x
      constructor

      · rintro ⟨⟨y, z⟩, ⟨hy, hz⟩, rfl⟩ j
        exact Fin.cases hy (fun k => hz k) j

      · intro hx
        refine ⟨(x 0, fun j => x j.succ), ⟨hx 0, fun j => hx j.succ⟩, ?_⟩
        funext j
        exact Fin.cases rfl (fun k => rfl) j

    rw [← heq]
    exact compact_image _ hk f hf
  \end{lean}

\end{proof}

% Lean source: Textbooks.MathematicalAnalysis1.Chapter01.CompactSets.continuous_to_euclidean
\begin{lemma} \label{an1:lem:continuous-to-euclidean}
  The coordinate identity from the finite-product metric on $\mathbb R^n$ to the Euclidean metric is continuous, including $n=0$.

  \begin{lean}
    theorem continuous_to_euclidean (n : ℕ) :
        Continuous (fun x : Fin n → ℝ => (WithLp.toLp 2 x : EuclideanSpace ℝ (Fin n)))
  \end{lean}
\end{lemma}

\begin{proof}
  Fix $\varepsilon>0$ and take $\delta=\varepsilon/(n+1)>0$. For two vectors within product distance $\delta$, let $D$ be that product distance and $d$ their Euclidean distance.

  \begin{lean}
    apply Metric.continuous_iff.mpr
    intro x ε hε

    have hn : (0 : ℝ) ≤ n := Nat.cast_nonneg n

    have hn1 : (0 : ℝ) < n + 1 := by
      positivity

    refine ⟨ε / (n + 1), div_pos hε hn1, ?_⟩
    intro y hy

    let D := dist y x

    let d := dist (WithLp.toLp 2 y : EuclideanSpace ℝ (Fin n)) (WithLp.toLp 2 x)

    have hD : 0 ≤ D := dist_nonneg

    have hd : 0 ≤ d := dist_nonneg
  \end{lean}

  Each coordinate distance is at most $D$. Squaring these nonnegative inequalities and summing the $n$ terms gives $d^2\le nD^2$ by the Euclidean distance formula.

  \begin{lean}
    have hsquare : d ^ 2 ≤ (n : ℝ) * D ^ 2 := by
      rw [PiLp.dist_sq_eq_of_L2]
      calc
        ∑ j : Fin n, dist (y j) (x j) ^ 2 ≤ ∑ _j : Fin n, D ^ 2 := by
          apply Finset.sum_le_sum
          intro j _

          have hle : dist (y j) (x j) ≤ D := dist_le_pi_dist y x j

          have hnonneg : 0 ≤ dist (y j) (x j) := dist_nonneg

          nlinarith
        _ = (n : ℝ) * D ^ 2 := by
          simp only [Finset.sum_const, Finset.card_univ,
            Fintype.card_fin, nsmul_eq_mul]
  \end{lean}

  Since $n\le(n+1)^2$ and both distances are nonnegative, $d\le(n+1)D$. The choice $D<\varepsilon/(n+1)$ therefore gives $d<\varepsilon$. This estimate also covers the empty coordinate type.

  \begin{lean}
    have hbound : d ≤ ((n : ℝ) + 1) * D := by
      have hextra : 0 ≤ ((n : ℝ) ^ 2 + n + 1) * D ^ 2 :=
        mul_nonneg (by positivity) (sq_nonneg D)

      apply (sq_le_sq₀ hd (mul_nonneg hn1.le hD)).mp
      nlinarith only [hsquare, hextra]

    have hsmall : ((n : ℝ) + 1) * D < ε := by
      have h := (lt_div_iff₀ hn1).mp hy

      nlinarith

    exact hbound.trans_lt hsmall
  \end{lean}

\end{proof}

% Lean source: Textbooks.MathematicalAnalysis1.Chapter01.CompactSets.cell_compact
\begin{theorem} \label{an1:thm:k-cell-is-compact}
  For every $n\in\N$ and real endpoints $a_j\leq b_j$, the cell $\prod_{j=0}^{n-1}[a_j,b_j]$ is compact in $\R^{n}$.

  \begin{lean}
    theorem cell_compact (n : ℕ) (a b : Fin n → ℝ) (hab : ∀ j, a j ≤ b j) :
        IsCompact {x : EuclideanSpace ℝ (Fin n) | ∀ j, a j ≤ x j ∧ x j ≤ b j}
  \end{lean}
\end{theorem}

\begin{proof}
  The preceding distance estimate proves continuity of the coordinate identity from the product description to the Euclidean description. Its image is exactly the same coordinate box, so the preceding result and compactness of continuous images apply. This includes $n=0$.

  \begin{lean}
    have h := compact_image _ (pi_cell_compact n a b hab)
      (fun x : Fin n → ℝ => (WithLp.toLp 2 x : EuclideanSpace ℝ (Fin n)))
      (continuous_to_euclidean n)

    convert h using 1
    ext x
    constructor

    · intro hx
      exact ⟨WithLp.ofLp x, hx, rfl⟩

    · rintro ⟨y, hy, rfl⟩
      exact hy
  \end{lean}

\end{proof}

% Lean source: Textbooks.MathematicalAnalysis1.Chapter01.CompactSets.closed_interval_compact
\begin{corollary} \label{an1:cor:closed-interval-in-R-is-compact}
  For $a,b\in\R$ with $a\leq b$, the closed interval $[a,b]$ is compact.

  \begin{lean}
    theorem closed_interval_compact (a b : ℝ) (hab : a ≤ b) : IsCompact (Icc a b)
  \end{lean}
\end{corollary}

\begin{proof}
  This is the real interval result established in the construction of compact cells.

  \begin{lean}
    exact interval_compact a b hab
  \end{lean}

\end{proof}

% Lean source: Textbooks.MathematicalAnalysis1.Chapter01.CompactSets.heine_borel
\begin{theorem} \label{an1:thm:heine-borel-theorem}
  For $n\in\N$ and $E\subseteq\R^{n}$, $E$ is compact if and only if it is closed and bounded. This is the \textbf{Heine--Borel theorem}.

  \begin{lean}
    theorem heine_borel (n : ℕ) (E : Set (EuclideanSpace ℝ (Fin n))) :
        IsCompact E ↔ IsClosed E ∧ Bornology.IsBounded E
  \end{lean}
\end{theorem}

\begin{proof}
  A compact set is closed and bounded by the metric-space results above. Conversely, suppose $E$ is closed and bounded. Euclidean space has a zero vector even in dimension zero, so choose a ball $B(p,r)$ containing $E$.

  \begin{lean}
    constructor

    · intro h
      exact ⟨compact_closed E h, compact_bounded E h⟩

    · rintro ⟨hc, hb⟩

      obtain ⟨p, r, hr, hs⟩ := (isBounded_iff E).mp hb ⟨0⟩
  \end{lean}

  The box with coordinate intervals $[p_j-r,p_j+r]$ is compact. Each coordinate difference is bounded in absolute value by the Euclidean norm, so a point in $B(p,r)$ lies in this box. Hence $E$ is a closed subset of a compact set and is compact.

  \begin{lean}
    let K : Set (EuclideanSpace ℝ (Fin n)) := {x | ∀ j, p j - r ≤ x j ∧ x j ≤ p j + r}

    have hk : IsCompact K := cell_compact n _ _ (fun j => by linarith)

    apply closed_subset_compact K E hk hc
    intro x hx j

    have hd : ‖x - p‖ < r := by
      simpa only [mem_ball, dist_eq_norm] using hs hx

    have hj : ‖(x - p) j‖ ≤ ‖x - p‖ := PiLp.norm_apply_le _ _

    change ‖x j - p j‖ ≤ ‖x - p‖ at hj
    rw [Real.norm_eq_abs] at hj

    have hcoord := abs_le.mp (hj.trans (le_of_lt hd))

    constructor <;> linarith
  \end{lean}

\end{proof}
